\centerline{\bf \Huge Пятая точка}

\begin{flushright}
        \scriptsize{}
\end{flushright}

\problem{1}
Три прямые пересекаются под углом $60^{\circ}$ в точке $P$. Из точки $O$, не лежащей на прямых, на них опущены перпендикуляры $O A, O B$ и $O C$. Докажите, что треугольник $A B C$ --- правильный.

\problem{2}
На плоскости дана окружность $\omega$ и точка $P$. Через точку $P$ проведены четыре прямые, каждая из которых пересекает $\omega$ в двух точках. Докажите, что середины четырех образовавшихся хорд лежат на одной окружности.

\problem{3}
В треугольнике $A B C$ провели высоты $A A^{\prime}, B B^{\prime}$ и $C C^{\prime}$.

(a) Точку $P$ спроектировали на высоты треугольника. Докажите, что треугольник, для которого эти три проекции являются вершинами, подобен треугольнику $A B C$.

(b) Пусть $B_1, C_1$ - середины высот $B B^{\prime}$ и $C C^{\prime}$. Докажите, что треугольники $A^{\prime} B_1 C_1$ и $A B C$ подобны.

(c) На высотах взяли точки $A_2, B_2, C_2$ так, что они делят высоты в отношении $2: 1$, считая от вершины. Докажите, что треугольники $A_2 B_2 C_2$ и $A B C$ подобны.

\problem{4}
Пусть $I$ - центр вписанной окружности треугольника $A B C$, а $I_a$ и $I_b$ ---1 центры вневписанных окружностей, касающихся сторон $B C$ и $A C$ соответственно. Докажите, что точки $I, C$ и середины отрезков $A I_a, B I_b$ лежат на одной окружности.

\problem{5}
Пусть $A_1$ - точка описанной окружности треугольника $A B C$, противоположная точке $A$. Точка $A_2$ симметрична точке $A_1$ относительно прямой $B C$. Аналогично определяются точки $B_2$ и $C_2$. Докажите, что описанная окружность треугольника $A_2 B_2 C_2$ проходит через точку $H$.